% Part: turing-machines
% Chapter: undecidability
% Section: verification

\documentclass[../../../include/open-logic-section]{subfiles}

\begin{document}

\olfileid{tur}{und}{ver}
\olsection{د تمثيل تاييدول}

\begin{explain}
د دې د تاييدولو دپاره چې زموږ تمثيل کار کوي، بايد دوه خبرې ثابتې
کړو۔ لومړے بايد وښيو چې که~$M$ پر ننوت~$w$ ودرېږي، نو~$!T(M, w)
\lif !E(M, w)$ معتبر دے۔ بيا بايد معکوسه خبره وښيو، يعنې که~$!T(M,
w) \lif !E(M, w)$ معتبر وي، نو~$M$ په رښتيا هم پر ننوت~$w$ د چلولو
پر مهال بالاخره درېږي۔

د دې دواړو د ثبوت تګلارې بېخي سره توپير لري۔ د لومړۍ پايلې دپاره
بايد وښيو چې د لومړۍ درجې منطق يو !!a{sentence} (يعنې~$!T(M, w)
\lif !E(M, w)$) معتبر دے۔ د دې تر ټولو اسانه لار د يوه
!!a{derivation} ورکول دي۔ خو زموږ ثبوت بايد د ټولو~$M$ او~$w$ دپاره
کار وکړي، نو په اصل کښې داسې يو واحد !!{sentence} نشته چې موږ ئې
دپاره !!a{derivation} ورکړو، بلکې نامتناهي ډېر دي۔ نو تر ټولو ښه کار
چې کولے ئې شو دا دے چې په استقرا ثابته کړو: $M$ او~$w$ هر ډول وي، او
پر ننوت~$w$ د~$M$ درېدل هر څومره ګامونه اخلي، د~$!T(M, w) \lif
!E(M, w)$ يو !!a{derivation} به موجود وي۔

په طبيعي ډول به زموږ استقرا د هغو ګامونو پر شمېر مخته ځي چې~$M$ ئې
درېدني تشکيل ته تر رسېدو مخکې اخلي۔ په خپل استقرايي ثبوت کښې به ټينګه
کړو چې پر ننوت~$w$ د~$M$ د چلولو د هر ګام~$n$ دپاره~$!T(M, w)
\Entails !C(M, w, n)$، چې~$!C(M, w, n)$ له~$n$ ګامونو وروسته پر~$w$
د چلېدلي~$M$ تشکيل په سمه توګه بيانوي۔ اوس که~$M$ پر ننوت~$w$، مثلاً
له~$n$ ګامونو وروسته، ودرېږي، $!C(M, w, n)$ به يو درېدنی تشکيل بيان
کړي۔ دا به هم وښيو چې هر کله~$!C(M, w, n)$ يو درېدنی تشکيل بيانوي،
$!C(M, w, n) \Entails !E(M, w)$۔ نو که~$M$ پر ننوت~$w$ ودرېږي، د
کوم~$n$ دپاره به~$M$ له~$n$ ګامونو وروسته په درېدني تشکيل کښې وي۔
له همدې امله~$!T(M, w) \Entails !C(M, w, n)$، چې~$!C(M, w, n)$ يو
درېدنی تشکيل بيانوي؛ او څرنګه چې په هغه حالت کښې~$!C(M, w, n)
\Entails !E(M, w)$، نو~$!T(M, w) \Entails !E(M, w)$ تر لاسه کوو؛
يعنې~$\Entails !T(M, w) \lif !E(M, w)$۔
\textbf{د ژباړې د نسخې سمون (OLCMP-053):} سرچينې د دې استدلال په
وروستۍ نتيجه کښې د تمثيلوونکې غونډلې د نښې لومړے سمبول پرې ايښے و؛
دلته بېرته راوړل شوے دے۔

د معکوسې خبرې تګلاره بېخي توپير لري۔ دلته فرضوو چې~$\Entails !T(M,
w) \lif !E(M, w)$، او بايد ثابته کړو چې~$M$ پر ننوت~$w$ درېږي۔ له
فرض څخه مو تر لاسه کېږي چې~$!T(M, w) \Entails !E(M, w)$، يعنې
$!E(M, w)$ په هر هغه !!{structure} کښې رښتيا دے چې~$!T(M, w)$ پکې
رښتيا دے۔ نو داسې يو !!a{structure}~$\Struct{M}$ بيان کوو چې~$!T(M,
w)$ پکې رښتيا دے: ډومېن به ئې~$\Nat$ وي، او د ټولو~$\Obj Q_q$ او
$\Obj S_\sigma$ تعبير به پر ننوت~$w$ د~$M$ د چلولو پر مهال د تشکيلونو
له خوا ورکول کېږي۔ نو مثلاً~$\Sat{M}{\Obj Q_q(\num{m}, \num{n})}$
هله او يوازې هله چې~$M$ پر ننوت~$w$ د~$n$ ګامونو دپاره وچلول شي،
په حالت~$q$ کښې وي او مربع~$m$ لولي۔ اوس څرنګه چې د فرض له مخې
$!T(M, w) \Entails !E(M, w)$، او د جوړښت له مخې~$\Sat{M}{!T(M, w)}$،
نو~$\Sat{M}{!E(M, w)}$۔ خو~$\Sat{M}{!E(M, w)}$ هله او يوازې هله چې
کوم~$n \in \Domain{M} = \Nat$ موجود وي داسې چې پر ننوت~$w$ چلېدلے
$M$ له~$n$ ګامونو وروسته په درېدني تشکيل کښې وي۔
\textbf{د ژباړې د نسخې سمون (OLCMP-054):} سرچينې د جوړښت د صدق په
تشريح کښې په تېروتنه د ماشين پر ځاے د تمثيلوونکې غونډلې نوم راوړے و؛
دلته فاعل د بحث له ماشين سره سم شوے دے۔
\end{explain}

\begin{defn}
$!C(M, w, n)$ دا !!{sentence} وي:
\[
\Obj Q_q(\num{m}, \num{n}) \land \Obj S_{\sigma_0}(\num{0}, \num{n})
\land \dots \land \Obj S_{\sigma_k}(\num{k}, \num{n}) \land
\lforall[x][(\num{k} < x \lif \Obj S_\TMblank(x, \num{n}))]
\]
چې~$q$ په وخت~$n$ کښې د~$M$ حالت دے، $M$ په وخت~$n$ کښې مربع~$m$
لولي، مربع~$i$ په وخت~$n$ کښې نښه~$\sigma_i$ لري چې~$0 \le i \le
k$، او~$k$ يا په وخت~$0$ کښې د پټې تر ټولو ښي ناتشه مربع ده، يا هغه
تر ټولو ښي مربع ده چې د پټې سر له~$n$ ګامونو وروسته ليدلې ده؛ په دې
دواړو کښې هره يوه چې لويه وي۔
\end{defn}

\begin{lem}\ollabel{lem:halt-config-implies-halt}
که پر ننوت~$w$ چلېدلے~$M$ له~$n$ ګامونو وروسته په يوه درېدني تشکيل
کښې وي، نو~$!C(M, w, n) \Entails !E(M, w)$۔
\end{lem}

\begin{proof}
فرض کړئ~$M$ د ننوت~$w$ دپاره له~$n$ ګامونو وروسته درېږي۔ داسې يو
حالت~$q$، مربع~$m$، او نښه~$\sigma$ شته چې:
\begin{enumerate}
\item له~$n$ ګامونو وروسته~$M$ په حالت~$q$ کښې دے، هغه مربع~$m$
  لولي چې~$\sigma$ پرې ښکاري۔
\item د حالت بدلون تابع~$\delta(q, \sigma)$ ناتعريفه ده۔
\end{enumerate}
$!C(M, w, n)$ د دې تشکيل بيان دے او دا فقرې~$\Obj Q_{q}(\num{m},
\num{n})$ او~$\Obj S_{\sigma}(\num{m}, \num{n})$ به پکې شاملې وي۔
دا فقرې په ګډه~$!E(M, w)$ راايجابوي:
\[
\lexists[x][\lexists[y][(\bigvee_{\tuple{q, \sigma} \in
      X}(\Obj Q_q(x, y) \land \Obj S_\sigma(x, y)))]]
\]
ځکه~$\Obj Q_{q}(\num{m}, \num{n}) \land \Obj S_{\sigma}(\num{m},
\num{n}) \Entails \bigvee_{\tuple{q, \sigma} \in X} (\Obj Q_q(\num{m},
\num{n}) \land \Obj S_{\sigma}(\num{m}, \num{n}))$، ځکه
$\tuple{q,\sigma} \in X$۔
\textbf{د ژباړې د نسخې سمون (OLCMP-055):} سرچينې په وروستي
استدلال کښې داسې لومړني حالت او نښه وکارول چې په فرض کښې نۀ وو ټاکل
شوي، او د نښې له محمول څخه ئې د ژبې نښه هم پرې ايښې وه؛ دلته هماغه
ټاکل شوي حالت او نښه او بشپړ محمول کارول شوي دي۔
\end{proof}

\begin{explain}
نو که~$M$ د ننوت~$w$ دپاره ودرېږي، داسې~$n$ شته چې~$!C(M, w, n)
\Entails !E(M,w)$۔ اوس به وښيو چې د هر وخت~$n$ دپاره~$!T(M, w)
\Entails !C(M, w, n)$۔
\end{explain}

\begin{lem}
\ollabel{lem:config}
د هر~$n$ دپاره، که~$M$ له~$n$ ګامونو وروسته لا نۀ وي درېدلے، نو
$!T(M, w) \Entails !C(M, w, n)$۔
\end{lem}

\begin{proof}
د استقرا بنسټ: که~$n = 0$ وي، د~$!C(M, w, 0)$ عاطفې د~$!T(M, w)$
هم عاطفې دي، نو د هغه له خوا ايجابېږي۔

استقرايي فرض: که~$M$ له~$n$-م ګام څخه مخکې نۀ وي درېدلے، نو
$!T(M,w) \Entails !C(M, w, n)$۔ بايد وښيو چې (تر هغه چې~$!C(M, w,
n)$ يو درېدنی تشکيل نۀ بيانوي) $!T(M, w) \Entails !C(M, w, n+1)$۔

فرض کړئ~$n \ge 0$ او پر~$w$ پيل شوي~$M$ له~$n$ ګامونو وروسته په
حالت~$q$ کښې مربع~$m$ لولي۔ څرنګه چې~$M$ له~$n$ ګامونو وروسته نۀ
درېږي، د~$M$ په پروګرام کښې بايد له لاندې درېو بڼو څخه يوه لارښوونه
وي:
\textbf{د ژباړې د نسخې سمون (OLCMP-056):} سرچينې استقرايي ګام يوازې
د مثبتو شاخصونو دپاره پيل کړے و او له صفرم تشکيل څخه لومړي تشکيل ته
بدلون ئې پرې ايښے و؛ دلته شرط صفر هم رانغاړي۔

\begin{enumerate}
\item \ollabel{right} $\delta(q, \sigma) = \tuple{q', \sigma', \TMright}$

\item \ollabel{left} $\delta(q, \sigma) = \tuple{q', \sigma', \TMleft}$

\item \ollabel{stay} $\delta(q, \sigma) = \tuple{q', \sigma', \TMstay}$
\end{enumerate}

موږ به دا درې واړه حالتونه په وار سره وڅېړو۔

\begin{enumerate}
\item فرض کړئ د~\olref{right} په بڼه يوه لارښوونه شته۔ د
  \olref[rep]{defn:tm-descr}\olref[rep]{rep-right} له مخې معنا ئې دا
  ده چې
\begin{align*}
& \lforall[x][\lforall[y][((\Obj Q_{q}(x,
  y) \land \Obj S_\sigma(x, y)) \lif {}]]\\
& \qquad (\Obj Q_{q'}(x',
  y') \land \Obj S_{\sigma'}(x, y') \land !A(x, y)))
\intertext{د~$!T(M,w)$ يوه عاطفه ده۔ دا لاندې !!{sentence} راايجابوي
  (کلي منځته راوړنه، $\num{m}$ د~$x$ دپاره او~$\num{n}$ د~$y$
  دپاره):}
& (\Obj Q_{q}(\num{m}, \num{n}) \land \Obj S_{\sigma}(\num{m},
\num{n})) \lif {}\\
& \qquad (\Obj Q_{q'}(\num{m}', \num{n}') \land
\Obj S_{\sigma'}(\num{m}, \num{n}') \land !A(\num{m}, \num{n})).
\intertext{د استقرايي فرض له مخې~$!T(M, w) \Entails !C(M, w, n)$؛
  يعنې:}
& \Obj Q_q(\num{m}, \num{n}) \land \Obj S_{\sigma_0}(\num{0}, \num{n})
\land \dots \land \Obj S_{\sigma_k}(\num{k}, \num{n}) \land \\
& \qquad\lforall[x][(\num{k} < x \lif \Obj S_\TMblank(x, \num{n}))]\\
\intertext{څرنګه چې له~$n$ ګامونو وروسته د پټې مربع~$m$ نښه~$\sigma$
  لري، اړونده عاطفه~$\Obj S_\sigma(\num{m}, \num{n})$ ده، نو دا
  راايجابوي:}
& \Obj Q_{q}(\num{m}, \num{n}) \land \Obj S_{\sigma}(\num{m},
   \num{n})
\intertext{اوس تر لاسه کوو:}
& \Obj Q_{q'}(\num{m}', \num{n}') \land \Obj S_{\sigma'}(\num{m},
  \num{n}') \land {}\\
& \qquad\Obj S_{\sigma_0}(\num{0}, \num{n}') \land \dots \land
  \Obj S_{\sigma_k}(\num{k}, \num{n}') \land {}\\
& \qquad \lforall[x][(\num{k} < x \lif \Obj S_\TMblank(x, \num{n}'))]
\end{align*}
په دې ډول: لومړۍ کرښه نېغ په نېغه د مخکيني شرطي د تالي څخه د مودوس
پونېنس په وسيله راځي۔ د منځنۍ کرښې هره عاطفه---چې
$S_{\sigma_m}(\num{m},\num{n}')$ ترې وتلې ده---په~$!C(M, w, n)$ کښې
له اړوندې عاطفې او~$!A(\num{m}, \num{n})$ څخه راځي۔

که~$m < k$ وي، $!T(M,w) \Entails \num{m} < \num {k}$
(\olref[rep]{prop:mlessk}) او د~$<$ د تعدي له مخې
$\lforall[x][(\num{k} < x \lif \num{m} < x)]$ لرو۔ که~$m = k$ وي،
نو~$\lforall[x][(\num{k} < x \lif \num{m} < x)]$ يوازې د منطق له
مخې راځي۔ بيا وروستۍ کرښه په~$!C(M, w, n)$ کښې له اړوندې عاطفې،
$\lforall[x][(\num{k} < x \lif \num{m} < x)]$، او~$!A(\num{m},
\num{n})$ څخه راځي۔ که~$m<k$ وي، دا لا همدا~$!C(M, w, n+1)$ دے۔

اوس فرض کړئ~$m=k$۔ په هغه حالت کښې د پټې سر له~$n+1$ ګامونو وروسته
مربع~$k+1$ هم ليدلې ده، چې اوس تر ټولو ښي ليدل شوې مربع ده۔ نو
$!C(M, w, n+1)$ يوه نوې عاطفه~$\Obj S_\TMblank(\num{k}',\num{n}')$
لري، او وروستۍ عاطفه ئې~$\lforall[x][(\num{k}' < x \lif \Obj
S_\TMblank(x, \num{n}'))]$ ده۔ بايد تاييد کړو چې دا دوه !!{sentence}s
هم ايجابېږي۔

موږ لا~$\lforall[x][(\num{k} < x \lif \Obj S_\TMblank(x,
  \num{n}'))]$ لرو۔ په ځانګړي ډول، له دې څخه~$\num{k} < \num{k}'
\lif \Obj S_\TMblank(\num{k}', \num{n}')$ تر لاسه کوو۔ له اکسيوم
$\lforall[x][x < x']$ څخه~$\num{k} < \num{k}'$ تر لاسه کوو۔ د مودوس
پونېنس له مخې~$\Obj S_\TMblank(\num{k}',\num{n}')$ راځي۔

همدارنګه، څرنګه چې~$!T(M,w) \Entails \num{k} < \num{k}'$، د~$<$
د تعدي اکسيوم موږ ته~$\lforall[x][(\num{k}' < x \lif \Obj
  S_\TMblank(x, \num{n}'))]$ راکوي۔ (د دې تاييد د تمرین په توګه
پرېږدو۔)
\textbf{د ژباړې د نسخې سمون (OLCMP-057):} سرچينې په دې دوو ځايونو
کښې نحوي ثابتېدنه کارولې وه، حال دا چې راجع شوې قضيه او د ټولې برخې
ادعا معنايي ايجاب بيانوي؛ دواړه اړيکې له همدې مخې يو شان شوې دي۔

\item فرض کړئ د~\olref{left} په بڼه يوه لارښوونه شته۔ بيا د
  \olref[rep]{defn:tm-descr}\olref[rep]{rep-left} له مخې،
\begin{align*}
& \lforall[x][\lforall[y][((\Obj Q_{q}(x', y) \land \Obj
    S_{\sigma}(x', y)) \lif {}]]\\
& \qquad (\Obj Q_{q'}(x, y') \land \Obj
  S_{\sigma'}(x', y') \land !A(x', y))) \land {}\\
& \lforall[y][((\Obj Q_{q}(\num{0}, y) \land \Obj S_{\sigma}(\num{0},
    y)) \lif {}]\\
& \qquad (\Obj Q_{q'}(\num{0}, y') \land \Obj S_{\sigma'}(\num{0}, y')
  \land !A(\num{0}, y)))
\intertext{د~$!T(M,w)$ يوه عاطفه ده۔ که~$m>0$ وي، نو~$l = m - 1$
  پرېږدئ (يعنې~$m = l+1$)۔ د پورتني !!{sentence} لومړۍ عاطفه دا
  راايجابوي:}
& (\Obj Q_{q}(\num{l}', \num{n}) \land \Obj S_{\sigma}(\num{l}', \num{n}))
\lif {} \\
& \qquad (\Obj Q_{q'}(\num{l}, \num{n}') \land \Obj S_{\sigma'}(\num{l}',
\num{n}') \land !A(\num{l}', \num{n}))
\intertext{که نۀ وي، $l = m = 0$ پرېږدئ او لاندې هغه !!{sentence}
  وګورئ چې دويمه عاطفه ئې راايجابوي:}
& ((\Obj Q_{q}(\num{0}, \num{n}) \land \Obj S_{\sigma}(\num{0}, \num{n})) \lif {}\\
& \qquad   (\Obj Q_{q'}(\num{0}, \num{n}') \land \Obj S_{\sigma'}(\num{0}, \num{n}') \land
!A(\num{0}, \num{n})))
\intertext{هره يوه غونډله دا راايجابوي:}
&  \Obj Q_{q'}(\num{l}, \num{n}') \land \Obj S_{\sigma'}(\num{m},
  \num{n}') \land {}\\
&\qquad  \Obj S_{\sigma_0}(\num{0}, \num{n}')\land \dots \land
  \Obj S_{\sigma_k}(\num{k}, \num{n}')  \land {} \\
& \qquad  \lforall[x][(\num{k} < x
  \lif \Obj S_\TMblank(x, \num{n}'))]
\end{align*}
لکه مخکې۔ (په ياد ولرئ چې په لومړي حالت کښې~$\num{l}' \ident
\num{l+1} \ident \num{m}$، او په دويم حالت کښې~$\num{l} \ident
\num{0}$۔) خو دا کټ مټ~$!C(M, w, n+1)$ دے۔
\textbf{د ژباړې د نسخې سمون (OLCMP-058):} سرچينې د کيڼ-حرکت په
لومړۍ څانګه کښې د ليکل شوې مربع پر ځاے ګاونډۍ مربع له بدلون څخه
مستثنا کړې وه، او د صفرمې مربع په څانګه کښې ئې د روانې لارښوونې پر
ځاے بې‌اړيکې شاخصي حالتونه کارولي وو۔ اکسيوم او دواړه منځته راوړنې
دلته د هماغې روانې لارښوونې له حالتونو او ليکل شوې مربع سره سمې شوې
دي۔

\item د~\olref{stay} حالت د تمرین په توګه پرېښوول شوے دے۔
\end{enumerate}
موږ وښودل چې د هر~$n$ دپاره~$!T(M, w) \Entails !C(M, w, n)$۔
\end{proof}

\begin{prob}
د~\olref[tur][und][ver]{lem:config} د ثبوت حالت
\olref[tur][und][ver]{stay} بشپړ کړئ۔
\end{prob}

\begin{prob}
له~$\Obj S_{\sigma_i}(\num{i}, \num{n})$ او~$!A(m, n)$ څخه د
$\Obj S_{\sigma_i}(\num{i}, \num{n}')$ يو !!a{derivation} ورکړئ
(په دې فرض چې~$i \neq m$؛ يعنې يا~$i < m$ يا~$m < i$)۔
\end{prob}

\begin{prob}
له~$\lforall[x][(\num{k} < x \lif \Obj S_\TMblank(x, \num{n}'))]$،
$\lforall[x][x < x']$، او
$\lforall[x][\lforall[y][\lforall[z][ ((x < y \land y < z) \lif x <
      z)]]]$ څخه د~$\lforall[x][(\num{k}' < x \lif \Obj
  S_\TMblank(x, \num{n}'))]$ يو !!a{derivation} ورکړئ۔)
\end{prob}


\begin{lem}
\ollabel{lem:valid-if-halt}
که~$M$ پر ننوت~$w$ ودرېږي، نو~$!T(M, w) \lif
!E(M, w)$ معتبر دے۔
\end{lem}

\begin{proof}
د~\olref{lem:config} له مخې پوهېږو چې د هر وخت~$n$ دپاره په
وخت~$n$ کښې د~$M$ د تشکيل بيان~$!C(M, w, n)$ د~$!T(M, w)$ له خوا
ايجابېږي۔ فرض کړئ~$M$ له~$k$ ګامونو وروسته درېږي۔ په هغه وخت کښې به
د کوم~$m \in \Nat$ دپاره مربع~$m$ لولي۔ بيا~$!C(M, w, k)$ د~$M$
يو درېدنی تشکيل بيانوي؛ يعنې~$\Obj Q_q(\num{m}, \num{k})$ او~$\Obj
S_\sigma(\num{m}, \num{k})$ دواړه د عاطفو په توګه لري، چې
$\delta(q,\sigma)$ ناتعريفه ده۔ نو د
\olref{lem:halt-config-implies-halt} له مخې~$!C(M, w, k) \Entails
!E(M, w)$۔ خو څرنګه چې~$!T(M, w) \Entails !C(M, w, k)$، نو~$!T(M,
w) \Entails !E(M, w)$، او ځکه~$!T(M, w) \lif !E(M, w)$ معتبر دے۔
\end{proof}


\begin{explain}
د خپلې ادعا د تاييد بشپړولو دپاره بايد معکوس لوری هم ټينګ کړو: که
$!T(M, w) \lif !E(M, w)$ معتبر وي، نو~$M$ په رښتيا هم پر ننوت~$w$
له پيل وروسته درېږي۔
\end{explain}

\begin{lem}
\ollabel{lem:halt-if-valid}
که~$\Entails !T(M, w) \lif !E(M, w)$، نو~$M$ پر ننوت~$w$ درېږي۔
\end{lem}

\begin{proof}
د~$\Lang L_M$ هغه !!{structure}~$\Struct M$ وګورئ چې ډومېن ئې
$\Nat$ دے، $\Obj 0$ په~$0$، $\prime$ په تالي تابع، او~$<$ د کوچنيوالي
په اړيکه تعبيروي، او محمولونه~$\Obj Q_q$ او~$\Obj S_\sigma$ په دې
ډول تعبيروي:
\begin{align*}
  \Assign{\Obj Q_q}{M} & =
\Setabs{\tuple{m, n}}{\begin{array}{ll}\text{پر~$w$ له پيل وروسته، له~$n$ ګامونو وروسته،}\\ \text{$M$ په حالت $q$
  کښې دے او مربع~$m$ لولي}
\end{array}} \\
\Assign{\Obj S_\sigma}{M} & = \Setabs{\tuple{m, n}}{\begin{array}{ll}
\text{پر~$w$ له پيل وروسته، له $n$ ګامونو وروسته،}\\ \text{مربع~$m$ د~$M$
  نښه~$\sigma$ لري}
\end{array}}
\end{align*}
په بله وينا، موږ !!{structure}~$\Struct{M}$ داسې جوړوو چې هغه څه ګام
په ګام بيان کړي چې پر ننوت~$w$ پيل شوے~$M$ په رښتيا کوي۔ په څرګند
ډول~$\Sat{M}{!T(M, w)}$۔ که~$\Entails !T(M, w) \lif !E(M, w)$، نو
$\Sat{M}{!E(M, w)}$ هم؛ يعنې:
\[
\Sat{M}{\lexists[x][\lexists[y][(\bigvee_{\tuple{q, \sigma} \in
      X}(\Obj Q_q(x, y) \land \Obj S_\sigma(x, y)))]]}.
\]
څرنګه چې~$\Domain{M} = \Nat$، بايد داسې~$m$، $n \in \Nat$ موجود
وي چې د کوم~$q$ او~$\sigma$ دپاره~$\Sat{M}{\Obj Q_q(\num{m},
\num{n}) \land \Obj S_\sigma(\num{m}, \num{n})}$، داسې چې
$\delta(q, \sigma)$ ناتعريفه وي۔ د~$\Struct M$ د تعريف له مخې، معنا
ئې دا ده چې پر ننوت~$w$ پيل شوے~$M$ له~$n$ ګامونو وروسته په
حالت~$q$ کښې دے او نښه~$\sigma$ لولي، او د حالت بدلون تابع ناتعريفه
ده؛ يعنې~$M$ درېدلے دے۔
\end{proof}

\end{document}
